\section{Conclusion}\label{sec:7}

This paper recasts the fuzzy DID design of \citet{fuzzy2018} as a debiased machine-learning problem. The Wald-DID and time-corrected estimands are rewritten as unconditional moments and augmented with closed-form first-step influence functions, yielding orthogonal scores whose Riesz representers are propensity-weighted cell indicators. DML-Wald and DML-TC inherit $\sqrt{n}$-consistency and asymptotic normality under a standard product-rate condition and achieve the semiparametric efficiency bound up to normalization. A data-driven one-sided trimming rule selects the threshold that minimizes a sample analog of the estimator's asymptotic variance relative to the squared first stage. The rule discards only observations whose probability of belonging to the treated group is close to one and preserves the first-stage information that symmetric rules waste.

Two lessons run through the simulations and the INPRES application. Orthogonality, not the choice of learner, drives the gains. Across Lasso, Ridge, elastic-net, and random forests DML-TC estimates agree within $0.05$ log points per year of schooling at both margins, while Monte Carlo bias falls by a factor of three to five relative to Lasso plug-ins. When the first-stage coefficient is small, the Wald ratio amplifies noise and DML-TC becomes the operational estimator. At the INPRES primary-completion margin DML-TC corrects an upward bias in the Lasso plug-in and tightens the cluster-robust confidence interval. Reporting both scores serves as a specification check. Agreement is reassuring, and disagreement signals that the conditional stability of treatment effects may fail.

The FSIF construction transfers to richer designs. Three extensions follow naturally: staggered adoption with multiple periods \citep{callaway2023ddml}, panel data with individual fixed effects, and district-level treatment with cluster-level cross-fitting. The cluster case is treated in Remark~\ref{rem:cluster} following \citet{chiang2022multiway} and \citet{belloni2017program}; the first two are left for future work. The one-sided optimality of the trimming rule is likewise stated as a conjecture that tracks the population tilting measure (Appendix~\ref{sec:app_trimming}); a formal treatment with the estimated measure is the natural next step.
